% Ad M. Brutum 1.3 --- headnote
% ad-brutum-01-03, 21 April 43 BC, Rome

\textit{Cicero to M. Junius Brutus, written from Rome ---
Perseus dateline \textit{Scr. Romae xi K. Mai., ut videtur,
a. 711 (43)}, i.e.\ 21 April 43 BC, the day after the news
of the victory at Mutina reached the city. The dateline
matches the meta date. The letter belongs to the brief
high-water moment of Cicero's second political life: the
consuls Hirtius and Pansa have driven Antony off from
Decimus Brutus, the boy Caesar (Octavian) has played his
agreed part, and Cicero on the day before the kalends-minus-
eleven (xii Kal. Mai., 20 April) has been escorted by an
enormous crowd up to the Capitol and placed on the Rostra
to receive the city's thanks. That triumphal day is the
substantive content of section 2. The reader should hold
the date in mind: Cicero does not yet know that Hirtius
fell at Mutina on 21 April, the very day of writing, or
that Pansa is dying of the wounds taken at Forum Gallorum.
Within a week, both consuls will be dead, the political
arithmetic that this letter assumes will have collapsed,
and the management of the young Caesar will be a problem
of a wholly different order.}

\textit{The letter is in three movements. Section 1 is
Cicero's report card on his three pupils --- the consuls
have turned out as he had described them, and Octavian's
\textit{indoles virtutis} is ``extraordinary,'' a judgement
Cicero will live to regret. The famous half-sentence
\textit{utinam tam facile eum florentem et honoribus et
gratia regere ac tenere possimus quam facile adhuc
tenuimus}, ``I only wish we may guide and hold him as
easily when he is in full flower of office and popularity
as we have held him so far,'' is one of the most candid
admissions of strategy in the correspondence: the boy is
a piece on the board, and the board is changing. Section
2 records the ovation of 20 April --- the only moment in
Cicero's career when a Roman crowd treated him as the
saviour Catiline's day had implied but never granted ---
and the equally Ciceronian gloss that ``there is nothing
empty in me, nor should there be.'' Section 3 turns to
Brutus: Cicero asks for news, warns him against being too
generous (a coded reference to Brutus's clemency toward
his prisoner C. Antonius, brother of Mark Antony), and
states bluntly that in his own view ``the cause of the
three brothers is one and the same'' --- Mark Antony,
Lucius Antony, and Gaius Antony all deserve every
penalty the law allows.}
